\documentclass[9pt]{extarticle}
\usepackage{fontspec}
\setmainfont{Leelawadee UI}
\usepackage[a4paper,top=0.7cm,bottom=0.7cm,left=0.7cm,right=0.7cm]{geometry}
\usepackage{multicol}
\setlength{\columnsep}{0.45cm}
\setlength{\columnseprule}{0.4pt}
\usepackage{amsmath,amssymb}
\usepackage{xcolor}
\usepackage{colortbl}
\usepackage{tikz}
\usepackage{pgfplots}
\pgfplotsset{compat=1.18}
\usepackage{enumitem}
\setlist{nosep,leftmargin=1.1em,itemsep=0pt,topsep=1pt}
\usepackage{titlesec}
\usepackage{tcolorbox}
\tcbuselibrary{skins}

\definecolor{navy}{RGB}{20,40,90}
\definecolor{gold}{RGB}{190,140,20}
\definecolor{lightbg}{RGB}{240,244,250}
\definecolor{solbg}{RGB}{247,247,240}
\definecolor{ideabg}{RGB}{236,246,238}
\definecolor{ideafr}{RGB}{60,120,70}

\titleformat{\section}{\normalfont\bfseries\normalsize\color{navy}}{\thesection}{0.4em}{}[{\color{gold}\titlerule[1pt]}]
\titlespacing*{\section}{0pt}{4pt}{2pt}
\titleformat{\subsection}{\normalfont\bfseries\color{gold}}{\thesubsection}{0.3em}{}
\titlespacing*{\subsection}{0pt}{2pt}{0pt}

\newtcolorbox{formulabox}{
  colback=lightbg, colframe=navy!60, boxrule=0.4pt,
  arc=1pt, left=3pt, right=3pt, top=1pt, bottom=1pt,
  fontupper=\footnotesize
}
\newtcolorbox{solbox}{
  colback=solbg, colframe=gold!70!black, boxrule=0.4pt,
  arc=1pt, left=3pt, right=3pt, top=1.5pt, bottom=1.5pt,
  fontupper=\footnotesize
}
\newtcolorbox{ideabox}{
  colback=ideabg, colframe=ideafr, boxrule=0.4pt,
  arc=1pt, left=3pt, right=3pt, top=1.5pt, bottom=1.5pt,
  fontupper=\footnotesize
}

\newcommand{\dist}[1]{\subsection*{\textcolor{navy}{\textbullet\ #1}}}

\setlength{\parindent}{0pt}
\setlength{\parskip}{0.5pt}

\pgfplotsset{
  every axis/.append style={
    label style={font=\tiny}, tick label style={font=\tiny},
    width=4.2cm, height=2.4cm, axis lines=left,
    tick align=outside, ymin=0,
    axis line style={-latex, thin}
  }
}
\setlength{\abovedisplayskip}{2pt}
\setlength{\belowdisplayskip}{2pt}
\setlength{\abovedisplayshortskip}{1pt}
\setlength{\belowdisplayshortskip}{1pt}

\pgfmathdeclarefunction{gauss}{3}{%
  \pgfmathparse{1/(#3*sqrt(2*pi))*exp(-((#1-#2)^2)/(2*#3^2))}%
}

\begin{document}
\begin{center}
{\Large\bfseries\color{navy} Chapter 3: การแจกแจงของตัวสถิติ (Sampling Distribution)}\\[2pt]
{\small สรุปเนื้อหา — อ.ดร.ธนวรรณ ประฮาดไชย \quad Department of Statistics, Faculty of Science, KKU}
\end{center}
\vspace{2pt}
{\color{gold}\hrule height 1.4pt}
\vspace{4pt}

\begin{multicols}{3}

\section*{1. ประชากร ตัวอย่าง และ CLT}
\textbf{ประชากร (Population)} คือกลุ่ม\allowbreak ทั้งหมดที่สนใจ\allowbreak ศึกษา \textbf{ตัวอย่าง (Sample)} คือส่วนหนึ่ง\allowbreak ที่สุ่มมาจาก\allowbreak ประชากร

\begin{formulabox}
\begin{tabular}{@{}l c c@{}}
 & \textbf{ประชากร} & \textbf{ตัวอย่าง}\\
 & (พารามิเตอร์) & (ตัวสถิติ)\\\hline
ค่าเฉลี่ย & $\mu$ & $\bar X$\\
ความแปรปรวน & $\sigma^2$ & $S^2$\\
สัดส่วน & $p$ & $\hat p$\\
\end{tabular}\\[2pt]
\footnotesize ตัวสถิติ (จากตัวอย่าง) ใช้\textbf{ประมาณค่า}หรือ\textbf{ทดสอบสมมติฐาน}เกี่ยวกับพารามิเตอร์ (ของประชากร)
\end{formulabox}

\textbf{การแจกแจงของตัวสถิติ (Sampling Distribution)} คือการแจกแจงความน่าจะเป็นของตัวสถิติ (เช่น $\bar X$) ที่ได้จากการสุ่มตัวอย่างซ้ำ ๆ

\dist{ทฤษฎีลิมิตสู่ส่วนกลาง (CLT)}
ถ้า $\bar X$ เป็นค่าเฉลี่ยของตัวอย่างขนาด $n$ ที่สุ่มจากประชากร\textbf{แบบใด ๆ} ที่มีค่าเฉลี่ย $\mu$ และความแปรปรวน $\sigma^2$ และ $n$ มีขนาด\textbf{ใหญ่พอ} ($n\ge30$) แล้ว
\begin{formulabox}
$$\bar X \sim N\left(\mu,\dfrac{\sigma^2}{n}\right) \text{ โดยประมาณ}$$
\end{formulabox}
\begin{ideabox}
\textbf{แนวคิด:} CLT สำคัญเพราะทำให้\allowbreak ใช้โค้งปกติ\allowbreak ประมาณค่าเฉลี่ยตัวอย่าง\allowbreak ได้ \textbf{โดยไม่ต้องรู้}\allowbreak ว่าประชากร\allowbreak แจกแจงแบบใด ขอเพียง $n$ มากพอ ($n\ge30$)
\end{ideabox}

กราฟด้านล่างแสดงว่าเมื่อ $n$ เพิ่มขึ้น การแจกแจงของ $\bar X$ จะยิ่งแคบและใกล้เคียงโค้งปกติมากขึ้น (ความแปรปรวนลดลงเป็น $\sigma^2/n$):
\begin{center}
\begin{tikzpicture}
\begin{axis}[width=6cm,height=3cm,axis lines=left,
  xlabel={$\bar x$},ylabel={},ymin=0,ymax=1.05,
  xtick={0},tick label style={font=\tiny},
  label style={font=\tiny},axis line style={-latex,thin},
  legend style={font=\tiny,draw=none,fill=none,at={(0.98,0.95)},anchor=north east}]
\addplot[domain=-4:4,samples=80,gray,thick] {gauss(x,0,1)/gauss(0,0,1)};
\addlegendentry{$n=1$}
\addplot[domain=-4:4,samples=80,gold,thick] {gauss(x,0,0.45)/gauss(0,0,0.45)};
\addlegendentry{$n=5$}
\addplot[domain=-4:4,samples=80,navy,thick] {gauss(x,0,0.2)/gauss(0,0,0.2)};
\addlegendentry{$n=30$}
\end{axis}
\end{tikzpicture}
\end{center}
{\footnotesize $n$ ยิ่งมาก โค้งของ $\bar X$ ยิ่งแคบ (แปรปรวนลดลง $=\sigma^2/n$) และเข้าใกล้โค้งปกติมากขึ้น\par}

\section*{2. การแจกแจงของค่าเฉลี่ยตัวอย่าง\\(Sampling distribution of $\bar X$)}

\dist{กรณีที่ 1: ทราบค่า $\sigma^2$}
ให้ $X_1,\dots,X_n$ เป็นตัวอย่างสุ่มขนาด $n$ จากประชากรที่มีการแจกแจง\textbf{แบบปกติ} ค่าเฉลี่ย $\mu$ ความแปรปรวน $\sigma^2$
\begin{formulabox}
$$\bar X \sim N\left(\mu,\dfrac{\sigma^2}{n}\right) \quad (n \text{ ขนาดใดก็ได้})$$
$$Z=\dfrac{\bar X-\mu}{\sigma/\sqrt{n}} \sim N(0,1)$$
\end{formulabox}
\begin{ideabox}
\textbf{แนวคิด:} รู้ $\sigma^2$ ของประชากร\allowbreak และประชากร\allowbreak แจกแจงปกติอยู่แล้ว $\to$ ใช้ $Z$ ได้ทันที ไม่ว่า $n$ จะเล็กหรือใหญ่ (ไม่ต้องพึ่ง CLT)
\end{ideabox}

\textit{ตัวอย่าง 1:} คะแนนสอบภาษาอังกฤษแจกแจงปกติ $\mu{=}500$, $\sigma{=}100$ สุ่มนักศึกษา $n{=}25$ คน จงหา ก) $P(\bar X{>}555)$ ข) $P(490{<}\bar X{<}515)$
\begin{solbox}
ก) $Z=\dfrac{555-500}{100/\sqrt{25}}=\dfrac{55}{20}=2.75$\\
$P(\bar X{>}555)=P(Z{>}2.75)$\\
$=0.5-P(0{<}Z{<}2.75)=0.5-0.4970=0.003$\\[3pt]
ข) $Z_1=\dfrac{490-500}{20}=-0.5,\ Z_2=\dfrac{515-500}{20}=0.75$\\
$P(490{<}\bar X{<}515)=P(-0.5{<}Z{<}0.75)$\\
$=P(0{<}Z{<}0.5)+P(0{<}Z{<}0.75)$\\
$=0.1915+0.2734=0.4649$
\end{solbox}

\dist{กรณีที่ 2: ไม่ทราบค่า $\sigma^2$}
ประมาณ $\sigma^2$ ด้วย $S^2=\dfrac{\sum(X_i-\bar X)^2}{n-1}$
\begin{formulabox}
ถ้า $n\ge30$ (อาศัย CLT):
$$Z=\dfrac{\bar X-\mu}{S/\sqrt{n}}\sim N(0,1) \text{ โดยประมาณ}$$
ถ้า $n<30$ (ต้องมี ประชากรปกติ):
$$T=\dfrac{\bar X-\mu}{S/\sqrt{n}}\sim t_{(n-1)}$$
\end{formulabox}
\begin{ideabox}
\textbf{แนวคิด:} ถ้า $n$ มากพอ ($\ge30$) ใช้ $S$ แทน $\sigma$ ได้และยัง\allowbreak ใช้ $Z$ ได้เพราะ\allowbreak CLT ช่วย แต่ถ้า $n$ น้อย ($<30$) การประมาณ\allowbreak $\sigma$ ด้วย $S$ ทำให้\allowbreak มีความแปรปรวนเพิ่ม จึงต้องใช้\allowbreak การแจกแจงแบบที (หางหนากว่าปกติ) แทน และ\textbf{ต้องอาศัยว่าประชากรปกติ}
\end{ideabox}

\dist{การแจกแจงแบบที (t-distribution)}
\begin{formulabox}
$$P(T_{(\nu)}>t_{(\alpha)})=\alpha,\quad -\infty<t<\infty$$
\end{formulabox}
รูปร่างคล้ายโค้งปกติ สมมาตรรอบ 0 แต่หางหนากว่า มีพารามิเตอร์เดียวคือ องศาความเป็นอิสระ $\nu$ (df); เมื่อ $\nu>30$ จะเข้าใกล้โค้งปกติมาตรฐาน

\begin{center}
\begin{tikzpicture}
\begin{axis}[width=4.6cm,height=2.4cm,axis lines=left,
  xlabel={},ylabel={},ymin=0,ymax=0.45,
  xtick={0},tick label style={font=\tiny},
  label style={font=\tiny},axis line style={-latex,thin},
  legend style={font=\tiny,draw=none,fill=none,at={(0.98,0.95)},anchor=north east}]
\addplot[domain=-4:4,samples=80,navy,thick] {gauss(x,0,1)};
\addlegendentry{$Z$ (ปกติ)}
\addplot[domain=-4:4,samples=80,gold,thick] {1/(sqrt(3*pi))*(1+x^2/3)^(-2)};
\addlegendentry{$t_{(3)}$}
\end{axis}
\end{tikzpicture}
\end{center}
{\footnotesize เส้นทอง ($t$, df น้อย) หางหนากว่าเส้นน้ำเงิน ($Z$)\par}

\textbf{ตารางที่ V} ให้ค่าวิกฤต $t_{(\alpha,\nu)}$ (พื้นที่หาง\textbf{ขวา} $=\alpha$) — อ่าน: หาแถว $\nu$(df) ตัดคอลัมน์ $\alpha$

\textit{ตัวอย่างอ่านตาราง:} $P(T_{(6)}>3.143)=0.01$ (แถว $\nu{=}6$ คอลัมน์ $\alpha{=}0.01$)

โดยสมมาตร: $P(T_{(\nu)}<-t_\alpha)=\alpha=P(T_{(\nu)}>t_\alpha)$

\textit{ตัวอย่าง:} $P(T_{(8)}<-1.860)=P(T_{(8)}>1.860)=0.05$

หรือใช้ Excel: \texttt{T.DIST}\allowbreak\texttt{(t,df,TRUE)} ให้ $P(T_{df}{<}t)$ (สะสมซ้าย)

\columnbreak
\textit{ตัวอย่าง 2:} ส่วนสูงนักเรียนอนุบาลแจกแจงปกติ $\mu{=}130$ ซม. สุ่ม $n{=}25$ คน ได้ $S{=}6$ ซม. จงหา $P(\bar X{>}131.1)$
\begin{ideabox}
\textbf{แนวคิด:} ไม่ทราบ $\sigma^2$ และ $n{=}25{<}30$ (แต่โจทย์บอกประชากร\allowbreak แจกแจงปกติ) $\to$ ต้องใช้การแจกแจง\allowbreak แบบที ไม่ใช่ $Z$
\end{ideabox}
\begin{solbox}
$T=\dfrac{131.1-130}{6/\sqrt{25}}=\dfrac{1.1}{1.2}=0.917$, $df{=}n{-}1{=}24$\\
จากตาราง V แถว $\nu{=}24$: $\alpha{=}.20\to t{=}0.857$; $\alpha{=}.10\to t{=}1.318$\\
เนื่องจาก $0.857<0.917<1.318$:\\
$0.10<P(T_{(24)}{>}0.917)<0.20$\\
\textit{(เอกสารต้นฉบับใช้ตารางที่มีคอลัมน์ $\alpha{=}.25$ แทน $.20$ จึงรายงานเป็น $0.10{<}P{<}0.25$)}
\end{solbox}

\textit{ตัวอย่าง 3:} รายได้ครัวเรือน\textbf{ไม่แจกแจงปกติ} $\mu{=}10{,}000$ บาท สุ่ม $n{=}100$ ครัวเรือน ได้ $S{=}4{,}800$ บาท จงหา $P(\bar X{<}10{,}500)$
\begin{ideabox}
\textbf{แนวคิด:} ประชากรไม่ปกติ และไม่ทราบ $\sigma^2$ แต่ $n{=}100\ge30$ $\to$ อาศัย \textbf{CLT} ใช้ $Z$ ได้ (แม้ประชากรไม่ปกติ เพราะ $n$ ใหญ่พอ)
\end{ideabox}
\begin{solbox}
$Z=\dfrac{10500-10000}{4800/\sqrt{100}}=\dfrac{500}{480}=1.04$\\
$P(\bar X{<}10500)=P(Z{<}1.04)$\\
$=0.5+P(0{<}Z{<}1.04)=0.5+0.3508=0.8508$
\end{solbox}

\section*{3. การแจกแจงของผลต่างค่าเฉลี่ย\\(Sampling distribution of $\bar X_1-\bar X_2$)}

\dist{กรณีที่ 1: ทราบ $\sigma_1^2,\sigma_2^2$}
\begin{formulabox}
$$\bar X_1-\bar X_2 \sim N\!\left(\mu_1-\mu_2,\ \dfrac{\sigma_1^2}{n_1}+\dfrac{\sigma_2^2}{n_2}\right)$$
$$Z=\dfrac{(\bar X_1-\bar X_2)-(\mu_1-\mu_2)}{\sqrt{\sigma_1^2/n_1+\sigma_2^2/n_2}}\sim N(0,1)$$
\end{formulabox}
\begin{ideabox}
\textbf{แนวคิด:} เปรียบเทียบค่าเฉลี่ย 2 ประชากร\allowbreak อิสระกัน เมื่อรู้ $\sigma_1^2,\sigma_2^2$ ทั้งคู่ $\to$ ใช้ $Z$ ตรง ๆ (ถ้า $n_1,n_2$ ไม่ใหญ่พอ ต้องอาศัย CLT เช่นกัน หากประชากรไม่ปกติ)
\end{ideabox}

\textit{ตัวอย่าง 1:} นิโคตินบุหรี่ A: $\mu_1{=}22,\sigma_1^2{=}36$; บุหรี่ B: $\mu_2{=}23,\sigma_2^2{=}64$ (มก.) สุ่มอย่างละ $n{=}100$ มวน จงหา $P$(นิโคติน A ต่ำกว่า B ไม่เกิน 0.5 มก.)
\begin{solbox}
ต้องการ $P(\bar X_1-\bar X_2<-0.5)$\\
$Z=\dfrac{-0.5-(22-23)}{\sqrt{36/100+64/100}}=\dfrac{0.5}{1}=0.5$\\
$P(\bar X_1{-}\bar X_2{<}{-0.5})=P(Z{<}0.5)$\\
$=0.5+P(0{<}Z{<}0.5)=0.5+0.1915=0.6915$
\end{solbox}

\dist{กรณีที่ 2: ไม่ทราบ $\sigma_1^2{=}\sigma_2^2$ (เท่ากัน)}
\begin{formulabox}
$$S_p^2=\dfrac{(n_1{-}1)S_1^2+(n_2{-}1)S_2^2}{n_1+n_2-2}$$
$$T=\dfrac{(\bar X_1-\bar X_2)-(\mu_1-\mu_2)}{\sqrt{S_p^2\left(\frac1{n_1}+\frac1{n_2}\right)}}\sim t_{(\nu)}$$
$$\nu=n_1+n_2-2$$
\end{formulabox}
\begin{ideabox}
\textbf{แนวคิด:} ไม่รู้ความแปรปรวน\allowbreak ประชากรทั้งคู่ แต่\allowbreak \textbf{สมมติว่าเท่ากัน} จึงรวม\allowbreak ข้อมูล 2 กลุ่มมาประมาณ\allowbreak ความแปรปรวนร่วม ($S_p^2$) แทนที่จะ\allowbreak ประมาณแยกกัน แล้วใช้การแจกแจง\allowbreak แบบที df$=n_1{+}n_2{-}2$
\end{ideabox}

\textit{ตัวอย่าง 2:} รายได้พนักงานบริษัท A,B แจกแจงปกติ $\mu_1{=}9{,}900$, $\mu_2{=}9{,}300$ บาท ไม่ทราบ $\sigma^2$ แต่เท่ากัน สุ่ม $n_1{=}15$ ($S_1{=}250$), $n_2{=}10$ ($S_2{=}200$) จงหา $P$(A มากกว่า B อย่างน้อย 450 บาท)
\begin{solbox}
$S_p^2=\dfrac{(14)(250)^2+(9)(200)^2}{23}=53695.65$, $\nu{=}23$\\
ต้องการ $P(\bar X_1-\bar X_2\ge450)$\\
$T=\dfrac{450-(9900-9300)}{\sqrt{53695.65(\frac1{15}+\frac1{10})}}=-1.586$\\
$=P(T_{(23)}\ge-1.586)=1-P(T_{(23)}{>}1.586)$
\end{solbox}
\begin{ideabox}
\textbf{แนวคิด (อ่านตารางเมื่อค่าไม่ตรงตาราง):} $1.586$ ไม่มีในตาราง V ที่ $\nu{=}23$ แต่มีค่าใกล้เคียง: $\alpha{=}.10\to t{=}1.319$ และ $\alpha{=}.05\to t{=}1.714$ เนื่องจาก $1.319<1.586<1.714$ จึง\textbf{ประกบค่า} $\alpha$ ไว้ระหว่างกลาง
\end{ideabox}
\begin{solbox}
$0.05<P(T_{(23)}{>}1.586)<0.10$\\
$\Rightarrow 0.90<1-P(T_{(23)}{>}1.586)<0.95$\\
ดังนั้น $0.90<P(\bar X_1-\bar X_2\ge450)<0.95$
\end{solbox}

\section*{4. การแจกแจงของสัดส่วนตัวอย่าง\\(Sampling distribution of $\hat p$)}
สุ่มตัวอย่างขนาด $n$ จากประชากร\allowbreak ที่มีสัดส่วน\allowbreak ลักษณะที่สนใจ $p$
\begin{formulabox}
$$\hat p \sim N\left(p,\dfrac{pq}{n}\right) \text{ โดยประมาณ},\quad q=1-p$$
$$Z=\dfrac{\hat p-p}{\sqrt{pq/n}}\sim N(0,1)$$
\end{formulabox}
\begin{ideabox}
\textbf{แนวคิด:} $\hat p$ คือสัดส่วนที่นับจาก\allowbreak ตัวอย่าง (คล้ายค่าเฉลี่ยของ\allowbreak ตัวแปร 0/1 แบบแบร์นูลลี) เมื่อ $n$ มากพอ CLT ทำให้\allowbreak ประมาณด้วยโค้งปกติได้ โดยมี\allowbreak ความแปรปรวน $pq/n$
\end{ideabox}

\textit{ตัวอย่าง 3:} นักศึกษาที่จะเรียนต่อปริญญาโท $p{=}0.3$ สุ่ม $n{=}60$ คน จงหา 1) $P(0.25{<}\hat p{<}0.35)$ 2) $P(\hat p{>}0.2)$
\begin{solbox}
$\sqrt{pq/n}=\sqrt{(0.3)(0.7)/60}=0.0592$\\[3pt]
1) $Z_1=\dfrac{0.25-0.3}{0.0592}=-0.85,\ Z_2=\dfrac{0.35-0.3}{0.0592}=0.85$\\
$P=P(-0.85{<}Z{<}0.85)=2(0.3023)=0.6046$\\[3pt]
2) $Z=\dfrac{0.2-0.3}{0.0592}=-1.69$\\
$P(\hat p{>}0.2)=P(Z{>}-1.69)$\\
$=0.5+P(0{<}Z{<}1.69)=0.5+0.4545=0.9545$
\end{solbox}

\section*{5. การแจกแจงของผลต่างสัดส่วน\\(Sampling distribution of $\hat p_1-\hat p_2$)}
สุ่มตัวอย่างขนาด $n_1,n_2$ อิสระกัน จาก 2 ประชากรสัดส่วน $p_1,p_2$
\begin{formulabox}
$$\hat p_1-\hat p_2 \sim N\!\left(p_1-p_2,\ \dfrac{p_1q_1}{n_1}+\dfrac{p_2q_2}{n_2}\right)$$
$$Z=\dfrac{(\hat p_1-\hat p_2)-(p_1-p_2)}{\sqrt{p_1q_1/n_1+p_2q_2/n_2}}\sim N(0,1)$$
\end{formulabox}
\begin{ideabox}
\textbf{แนวคิด:} เปรียบเทียบสัดส่วน 2 กลุ่มอิสระกัน (คล้ายหัวข้อ 3 แต่ใช้กับสัดส่วนแทนค่าเฉลี่ย) ไม่ต้องสมมติ $p_1{=}p_2$ เพราะสูตรแปรปรวนคำนวณแยกแต่ละกลุ่มอยู่แล้ว
\end{ideabox}

\textit{ตัวอย่าง 4:} กลุ่มรับวิตามินซี $p_1{=}0.15$ เป็นหวัด; กลุ่มไม่รับ $p_2{=}0.20$ เป็นหวัด สุ่มกลุ่มละ $n_1{=}n_2{=}50$ คน จงหา $P$(สัดส่วนกลุ่มรับวิตามินซีมากกว่ากลุ่มไม่รับ น้อยกว่า 0.05)
\begin{solbox}
$\sqrt{\frac{p_1q_1}{n_1}{+}\frac{p_2q_2}{n_2}}=\sqrt{\dfrac{(.15)(.85)}{50}{+}\dfrac{(.2)(.8)}{50}}=0.0758$\\
$Z=\dfrac{0.05-(0.15-0.2)}{0.0758}=\dfrac{0.1}{0.0758}=1.32$\\
$P(\hat p_1-\hat p_2{<}0.05)=P(Z{<}1.32)$\\
$=0.5+P(0{<}Z{<}1.32)=0.5+0.4066=0.9066$
\end{solbox}

\section*{6. ตารางอ้างอิงฟังก์ชัน Excel}
\begin{formulabox}
\begin{tabular}{@{}ll@{}}
\textbf{ที (t)} & \texttt{T.DIST}\\
& \texttt{(t,df,TRUE)}\\
\end{tabular}\\[2pt]
\footnotesize ให้ $P(T_{df}<t)$ (สะสมซ้าย) — ต่างจากตาราง V ที่พิมพ์ในหนังสือซึ่งให้ค่าวิกฤต\textbf{หางขวา}
\end{formulabox}

\textit{หมายเหตุ:} เอกสารนี้\allowbreak ครอบคลุมเฉพาะหัวข้อที่มีสอนใน\allowbreak สไลด์จริง (ค่าเฉลี่ย, ผลต่างค่าเฉลี่ย, สัดส่วน, ผลต่างสัดส่วน) — หัวข้อ\allowbreak ``การแจกแจงของความแปรปรวน'' และ\allowbreak ``อัตราส่วนความแปรปรวน'' ที่ระบุใน\allowbreak สารบัญ (สไลด์ 2) เป็นสีเทา/\allowbreak ไม่ได้สอนจริงในไฟล์นี้ จึง\textbf{ไม่รวมอยู่ในสรุปนี้}

\end{multicols}
\newpage
\section*{7. ภาคผนวก: ตารางที่ V (ค่าจริง)}

\subsection*{\textcolor{ideafr}{ตารางนับจากฝั่งไหน?}}
\begin{center}
\begin{minipage}{9cm}
\centering
\textbf{ตาราง V (การแจกแจงแบบที) ต่างจากตาราง I--III ในเอกสาร CH2!}\\
ตาราง V ให้ค่าวิกฤต $t_{(\alpha,\nu)}$ โดยนับพื้นที่หาง\textbf{ขวา} (ไม่ใช่สะสมจากซ้ายแบบตาราง I--III)
\begin{tikzpicture}
\begin{axis}[width=8cm,height=4cm,axis lines=left,
  xlabel={$t$},ylabel={$f(t)$},ymin=0,ymax=0.45,
  xtick={-3,-2,-1,0,1,2,3},tick label style={font=\small},
  label style={font=\small},axis line style={-latex,thin}]
\addplot[domain=-3.4:3.4,samples=80,navy,thick] {gauss(x,0,1)};
\addplot[domain=1:3.4,samples=60,fill=gold!55,draw=none] {gauss(x,0,1)} \closedcycle;
\draw[-latex,thick,gold!70!black] (axis cs:3.2,-0.08) -- (axis cs:1,-0.08)
  node[midway,below,font=\small,black]{นับจากขวา $\to$ ซ้าย};
\draw[dashed] (axis cs:1,0) -- (axis cs:1,0.24) node[above,font=\small]{$t_{(\alpha,\nu)}$};
\end{axis}
\end{tikzpicture}\\
{\footnotesize พื้นที่แรเงา (ขวาของ $t_{(\alpha,\nu)}$ ถึง $+\infty$) $=\alpha=P(T_\nu{>}t_{(\alpha,\nu)})$; ตรงข้ามกับตาราง Z ที่นับพื้นที่ทางซ้าย\par}
\end{minipage}
\end{center}

\textbf{วิธีอ่าน:} หาแถว $\nu$ (df, ซ้ายสุด) ตัดคอลัมน์ $\alpha$ (หัวตาราง) — ค่าตรงจุดตัดคือ $t_{(\alpha,\nu)}$ ที่ทำให้ $P(T_\nu>t_{(\alpha,\nu)})=\alpha$

\textbf{ถ้าค่า $t$ ที่คำนวณได้ไม่ตรงกับตาราง} (เช่นตัวอย่าง 2 ในหัวข้อ 3): ให้หาสองคอลัมน์ $\alpha$ ที่ใกล้เคียงที่สุด (ค่าน้อยกว่าและมากกว่า) แล้ว\textbf{ประกบค่า} $\alpha$ ไว้ระหว่างกลาง

\subsection*{\textcolor{navy}{ตาราง V: ค่าวิกฤตการแจกแจงแบบที $t_{(\alpha,\nu)}$}}
\begin{center}
{\small
\begin{tabular}{r|ccccccc}
\rowcolor{lightbg} $\nu\backslash\alpha$ & .20 & .10 & .05 & .025 & .01 & .005 & .001\\\hline
1 & 1.376 & 3.078 & 6.314 & 12.706 & 31.821 & 63.657 & 318.309\\
2 & 1.061 & 1.886 & 2.920 & 4.303 & 6.965 & 9.925 & 22.327\\
3 & 0.978 & 1.638 & 2.353 & 3.182 & 4.541 & 5.841 & 10.215\\
4 & 0.941 & 1.533 & 2.132 & 2.776 & 3.747 & 4.604 & 7.173\\
5 & 0.920 & 1.476 & 2.015 & 2.571 & 3.365 & 4.032 & 5.893\\
6 & 0.906 & 1.440 & \textbf{1.943} & 2.447 & \textbf{3.143} & 3.707 & 5.208\\
7 & 0.896 & 1.415 & 1.895 & 2.365 & 2.998 & 3.500 & 4.785\\
8 & 0.889 & 1.397 & \textbf{1.860} & 2.306 & 2.897 & 3.355 & 4.501\\
9 & 0.883 & 1.383 & 1.833 & 2.262 & 2.821 & 3.250 & 4.297\\
10 & 0.879 & 1.372 & 1.812 & 2.228 & 2.764 & 3.169 & 4.144\\
15 & 0.866 & 1.341 & 1.753 & 2.131 & 2.602 & 2.947 & 3.733\\
20 & 0.860 & 1.325 & 1.725 & 2.086 & 2.528 & 2.845 & 3.552\\
21 & 0.859 & 1.323 & 1.721 & 2.080 & 2.518 & 2.831 & 3.527\\
22 & 0.858 & 1.321 & 1.717 & 2.074 & 2.508 & 2.819 & 3.505\\
23 & 0.858 & \textbf{1.319} & \textbf{1.714} & 2.069 & 2.500 & 2.807 & 3.485\\
24 & \textbf{0.857} & \textbf{1.318} & 1.711 & 2.064 & 2.492 & 2.797 & 3.467\\
25 & 0.856 & 1.316 & 1.708 & 2.060 & 2.485 & 2.787 & 3.450\\
26 & 0.856 & 1.315 & 1.706 & 2.056 & 2.479 & 2.779 & 3.435\\
27 & 0.855 & 1.314 & 1.703 & 2.052 & 2.473 & 2.771 & 3.421\\
28 & 0.855 & 1.313 & 1.701 & 2.048 & 2.467 & 2.763 & 3.408\\
29 & 0.854 & 1.311 & 1.699 & 2.045 & 2.462 & 2.756 & 3.396\\
30 & 0.854 & 1.310 & 1.697 & 2.042 & 2.457 & 2.750 & 3.385\\
\end{tabular}
\par}
\end{center}
{\small ตัวเลขที่ \textbf{ตัวหนา} คือค่าที่ใช้ในตัวอย่างต่าง ๆ ของเอกสารนี้ (เช่น $\nu{=}6,t{=}3.143\to\alpha{=}.01$; $\nu{=}8,t{=}1.860\to\alpha{=}.05$; $\nu{=}24$: $\alpha{=}.20\to.857$, $\alpha{=}.10\to1.318$; $\nu{=}23$: $\alpha{=}.10\to1.319$, $\alpha{=}.05\to1.714$)\\[3pt]
เมื่อ $\nu>30$ ค่าวิกฤตของตาราง V จะเข้าใกล้ค่าของตาราง Z (ปกติมาตรฐาน) มากขึ้นเรื่อย ๆ}

\end{document}
